\documentclass[../../../main]{subfiles}

\ifSubfilesClassLoaded{%
    \zexternaldocument*{../../1.Category-Theory/exported/aux/main}%
    \zexternaldocument*{1.Additive-categories}
}{}

\begin{document}
\section{Abelian categories}

\begin{definition}[Abelian category] \label{def:abelian-category}
  An \hyperref[def:additive-category]{additive category} \(\category{A}\) is \textbf{abelian} if it verifies
  \begin{enumerate}
    \item[A.4)]\label{axiom:abelian-existence} \hyperref[def:categorical-kernel]{Kernels} and \hyperref[def:categorical-cokernel]{cokernels} exist in \(\category{A}\).
    \item[A.5)]\label{axiom:abelian-morphisms} Every \hyperref[def:mono-epi-morphism]{monomorphism} is the \hyperref[def:categorical-kernel]{kernel} of some morphism and every \hyperref[def:mono-epi-morphism]{epimorphism} is the \hyperref[def:categorical-cokernel]{cokernel} of some morphism.
  \end{enumerate}
\end{definition}

\begin{lemma} \label{lemma:kernel-of-cokernel-abelian}
  In an \hyperref[def:abelian-category]{abelian category} \(\category{A}\) every \hyperref[def:categorical-kernel]{kernel} is the \hyperref[def:categorical-kernel]{kernel} of its \hyperref[def:categorical-cokernel]{cokernel} and every \hyperref[def:categorical-cokernel]{cokernel} is the \hyperref[def:categorical-cokernel]{cokernel} of its \hyperref[def:categorical-kernel]{kernel}.
  In formulas
  \begin{align*}
    \ker (\coker \varphi) = \ker \varphi \qquad \coker (\ker \varphi) = \coker \varphi
  \end{align*}
\end{lemma}
\begin{proof}
  Lets check \hyperref[def:categorical-cokernel]{cokernels} this time.

  Let \(\coker \varphi: A \to B\) be the \hyperref[def:categorical-cokernel]{cokernel} of a morphism \(\varphi: Z \to A\).
  Since \(\category{A}\) is \hyperref[def:abelian-category]{abelian} \(\varphi\) has a kernel \(\ker \varphi: K \to A\).
  Note that \(Z \to A \to B\) is 0, therefore by the \hyperref[def:categorical-kernel]{kernel} universal property we have
  \begin{center}
    \begin{tikzcd}
      Z \arrow[rr, "\varphi"] \arrow[dr, dotted] & & A \arrow[r, "\coker \varphi"'] & B \\
      & K \arrow[ur, "\ker \varphi"']
    \end{tikzcd}
  \end{center}

  Consider a morphism \(A \to C\) such that \((K \to A \to C) = 0\).
  Now using the \hyperref[def:categorical-cokernel]{cokernel} universal property we obtain
  \begin{center}
    \begin{tikzcd}
      &&& C & \\
      Z \arrow[rr, "\varphi"] \arrow[dr, dotted] \arrow[urrr, bend left, dotted, "0"] & & A \arrow[rr, "\coker \varphi"'] \arrow[ur] & & \arrow[ul, dotted] B \\
      & K \arrow[ur, "\ker \varphi"']
    \end{tikzcd}
  \end{center}
  Therefore \(\coker \varphi\) satisfies the universal property of the \hyperref[def:categorical-cokernel]{cokernel} of \(\ker \varphi: K \to A\).
\end{proof}

\begin{lemma} \label{lemma:mono-and-epi-is-iso}
  Let \(\varphi: A \to B\) be a morphism in an \hyperref[def:abelian-category]{abelian category} \(\category{A}\).
  If \(\varphi\) is both a \hyperref[def:mono-epi-morphism]{monomorphism and epimorphism} then \(\varphi\) is an \hyperref[def:categorical-isomorphism]{isomorphism}.
\end{lemma}
\begin{proof}
  By \autoref{lemma:partial-sequence-additive} the \hyperref[def:categorical-kernel]{kernel} of \(\varphi\) is \(0 \to A\) and its \hyperref[def:categorical-cokernel]{cokernel} is \(B \to 0\).
  Furthermore by \autoref{lemma:kernel-of-cokernel-abelian} \(\varphi\) is the \hyperref[def:categorical-cokernel]{cokernel} of \(0 \to A\) and the \hyperref[def:categorical-kernel]{kernel} of \(B \to 0\).

  Consider the following diagram
  \begin{center}
    \begin{tikzcd}
        &   & B \arrow[d, "\id"] &   \\
      0 \arrow[r] & A \arrow[r, tail, two heads, "\varphi"] & B \arrow[r] & 0
    \end{tikzcd}
  \end{center}
  Since \((B \to B \to 0) = 0\) and \(\varphi\) is the \hyperref[def:categorical-kernel]{kernel} of \(B \to 0\) we have
  \begin{center}
    \begin{tikzcd}
        &   & B \arrow[d, "\id"] \arrow[dl, dotted, "\psi"'] &   \\
      0 \arrow[r] & A \arrow[r, tail, two heads, "\varphi"] & B \arrow[r] & 0
    \end{tikzcd}
  \end{center}
  Therefore since \(\varphi \psi = \id\) we have shown that \(\varphi\) has a right-inverse.

  Consider now the diagram
  \begin{center}
    \begin{tikzcd}
                  &   A & &   \\
      0 \arrow[r] & A \arrow[u, "\id"] \arrow[r, tail, two heads, "\varphi"] & B \arrow[r] & 0
    \end{tikzcd}
  \end{center}
  Since \((0 \to A \to A) = 0\) and \(\varphi\) is the \hyperref[def:categorical-cokernel]{cokernel} of \(0 \to A\) we have
  \begin{center}
    \begin{tikzcd}
                  &   A & &   \\
      0 \arrow[r] & A \arrow[u, "\id"] \arrow[r, tail, two heads, "\varphi"] & B \arrow[ul, dotted, "\eta"'] \arrow[r] & 0
    \end{tikzcd}
  \end{center}
  Therefore since \(\eta \varphi = \id\) we have shown that \(\varphi\) has a left-inverse.
\end{proof}

\subsection*{The factorization property}

\begin{theorem}[Factorization property] \label{thm:factorization-property}
  Every morphism \(\varphi: A \to B\) in an \hyperref[def:abelian-category]{abelian category} \(\category{A}\) can be decomposed as
  \begin{center}
    \begin{tikzcd}
      A \arrow[d, two heads, "\coker \ker \varphi"'] \arrow[r, "\varphi"] & B \\
      C \arrow[r, dotted, "\sim", "\tilde{\varphi}"'] & K \arrow[u, tail, "\ker \coker \varphi"']
    \end{tikzcd}
  \end{center}
\end{theorem}

Before proving this statement, we need to consider some auxiliary lemmas.
\begin{lemma} \label{lemma:image-coimage-properties}
  Let \(\varphi: A \to B\) be a morphism in an \hyperref[def:abelian-category]{abelian category} and let \(\iota: K \to B\) be the \hyperref[def:categorical-kernel]{kernel} of its \hyperref[def:categorical-cokernel]{cokernel}. Then
  \begin{itemize}
    \item \(\iota\) is a \hyperref[def:mono-epi-morphism]{monomorphism}.
    \item \(\varphi\) factors through \(\iota\).
    \item \(\iota\) is \hyperref[def:initial-terminal-object]{initial} with these properties.
  \end{itemize}
\end{lemma}
\begin{proof}~
  \begin{itemize}
    \item By \autoref{lemma:kernels-cokernels-are-mono-epi} \(\iota\) is a \hyperref[def:mono-epi-morphism]{monomorphism}.
    \item Since \(A \to B \to \coker \varphi\) is the zero morphism by the universal property of \hyperref[def:categorical-kernel]{kernels} there is an induced morphism \(K \to A\).
    \item We want to see that if \(\lambda: L \monoarrow B\) is a \hyperref[def:mono-epi-morphism]{monomorphism} that factors through \(\varphi\) then \(\iota\) factors uniquely through \(\lambda\)
          \begin{center}
            \begin{tikzcd}
              & L \arrow[dr, tail, "\lambda"] &  \\
              A \arrow[ur] \arrow[r] \arrow[rr, bend right, "\varphi"']  & K \arrow[u, dotted] \arrow[r, tail, "\iota"] & B
            \end{tikzcd}
          \end{center}

          Consider \(\coker \lambda\)
          \begin{center}
            \begin{tikzcd}
              & L \arrow[dr, tail, "\lambda"] &  \\
              A \arrow[ur] \arrow[r]  & K \arrow[r, tail, "\iota"] & B \arrow[dr, "\coker \lambda"'] \\
              &&& \operatorname{Cok}_{\lambda}
            \end{tikzcd}
          \end{center}

          Since \(\varphi\) factors through \(\lambda\) we have that \(A \to B \to \operatorname{Cok}_{\lambda} = 0\). Therefore by the universal property of the \hyperref[def:categorical-cokernel]{cokernel} we have an induced map
           \begin{center}
            \begin{tikzcd}
              & L \arrow[dr, tail, "\lambda"] &  \\
              A \arrow[ur] \arrow[r]  & K \arrow[r, tail, "\iota"] & B \arrow[dr, "\coker \lambda"'] \arrow[r, "\coker \varphi"] & \operatorname{Cok}_{\varphi} \arrow[d, dotted] \\
              &&& \operatorname{Cok}_{\lambda}
            \end{tikzcd}
          \end{center}

          Therefore since \(K \to \operatorname{Cok}_{\varphi} = 0\) we have that \(K \to B \xrightarrow{\coker \lambda} \operatorname{Cok}_{\lambda} = 0\).

          Finally since by \autoref{lemma:kernel-of-cokernel-abelian} \(\lambda: L \to B\) is the \hyperref[def:categorical-kernel]{kernel} of its \hyperref[def:categorical-cokernel]{cokernel} there exists a unique morphism \(K \to L\) making the diagram commute.
  \end{itemize}
\end{proof}

\begin{definition} \label{def:image-coimage}
  Let \(\varphi: A \to B\) be a morphism in an \hyperref[def:abelian-category]{abelian category}.
  The \textbf{image} of \(\varphi\) is
  \begin{align*}
    \im \varphi = \ker \coker \varphi
  \end{align*}
  and the \textbf{coimage} of \(\varphi\) is
  \begin{align*}
    \coim \varphi = \coker \ker \varphi
  \end{align*}
\end{definition}

\begin{notation}
  If \(\varphi: A \monoarrow B\) is a \hyperref[def:mono-epi-morphism]{monomorphism} we denote by \(\quot{B}{A}\) the image of \(\coker \varphi\).

  In this way the target of \(\coim \varphi\) is \(\quot{B}{\ker \varphi}\).
\end{notation}

\begin{lemma}
  Let \(\varphi: A \to B\) be a morphism in an \hyperref[def:abelian-category]{abelian category} with \(\im \varphi: K \to B\) and \(\coim \varphi: A \to C\).
  Then the induced morphisms \(A \to K\) and \(B \to C\) are, respectively, an \hyperref[def:mono-epi-morphism]{epimorphism} and a \hyperref[def:mono-epi-morphism]{monomorphism}.
\end{lemma}
\begin{proof}
  We will see that \(\overline{\varphi}: A \to K\) is a \hyperref[def:mono-epi-morphism]{epimorphism}. The other statement follows analogously.

  Consider the \hyperref[def:image-coimage]{image} of \(\overline{\varphi}\)
  \begin{center}
    \begin{tikzcd}
      K' \arrow[r, tail, "\im \overline{\varphi}"] & K \arrow[r, two heads, "\coker \overline{\varphi}"] & \operatorname{Cok}
    \end{tikzcd}
  \end{center}
  Since by \autoref{lemma:kernel-of-cokernel-abelian} every \hyperref[def:mono-epi-morphism]{epimorphism} is the \hyperref[def:categorical-cokernel]{cokernel} of its \hyperref[def:categorical-kernel]{kernel} we have that \(\coker \overline{\varphi} = \coker \im \overline{\varphi}\).
  Moreover since \(K' \monoarrow K\) is a \hyperref[def:mono-epi-morphism]{monomorphism} through which \(\overline{\varphi}\) factors we have
   \begin{center}
     \begin{tikzcd}
       A \arrow[r, "\varphi"] \arrow[d] \arrow[dr, "\overline{\varphi}"'] & B \\
       K \arrow[r, tail, "\im \overline{\varphi}"'] & K' \arrow[u, tail]
    \end{tikzcd}
  \end{center}
  Since \(\im \varphi\) is initial among this kind of morphisms by \autoref{lemma:image-coimage-properties} necessarily \(\im \overline{\varphi}\) is an \hyperref[def:categorical-isomorphism]{isomorphism}.

  Therefore \(0 = \coker \im \overline{\varphi} = \coker \overline{\varphi}\), proving that \(\overline{\varphi}\) is an \hyperref[def:mono-epi-morphism]{epimorphism} by \autoref{lemma:partial-sequence-additive}.
\end{proof}

To summarize, we have the following situation
\begin{center}
  \begin{tikzcd}
    & K \arrow[dr, tail, "\im \varphi"] \\
    A \arrow[dr, two heads, "\coim \varphi"'] \arrow[ur, two heads, "\overline{\varphi}"] \arrow[rr, "\varphi"] & & B \\
    & C \arrow[ur, tail, "\underline{\varphi}"']
  \end{tikzcd}
\end{center}

\begin{proof}[Proof of \autoref{thm:factorization-property}]
  By the universal property of \(\im \varphi\) and \(\coim \varphi\) there exists a unique morphism \(\psi: K \to C\) making the above diagram commute
  \begin{center}
    \begin{tikzcd}
      & K \arrow[dr, tail, "\im \varphi"] \arrow[dd, dotted, "\psi" near end] \\
      A \arrow[dr, two heads, "\coim \varphi"'] \arrow[ur, two heads, "\overline{\varphi}"] \arrow[rr, "\varphi" near start] & & B \\
      & C \arrow[ur, tail, "\underline{\varphi}"']
    \end{tikzcd}
  \end{center}

  Since \(\underline{\varphi} \circ \psi = \im \varphi\) is a \hyperref[def:mono-epi-morphism]{monomorphism} so is \(\psi\).
  Since \(\psi \circ \overline{\varphi} = \coim \varphi\) is a \hyperref[def:mono-epi-morphism]{epimorphism} so is \(\psi\).
  Therefore by \autoref{lemma:mono-and-epi-is-iso} \(\psi\) is an \hyperref[def:categorical-isomorphism]{isomorphism}.

  Taking \(\tilde{\varphi} = \psi^{-1}: C \to K\) we are done.
\end{proof}

\end{document}
